\chapter{LLMOps para Prompts e Sistemas de Produção}
\label{ch:llmops}

\ficcao A Lavandaria encurtou o prazo de devolução de dez para cinco dias úteis. A política foi publicada. O sistema continuou a responder «dez». Tinha cumprido todos os testes de formato, incluindo aquele que verificava se o prazo era um número. Era, de facto, um número.

Se a resposta saiu num chat de experiências, podemos corrigir o pedido e repetir. Num serviço em produção, precisamos de saber quem recebeu a resposta, que fonte a sustentou e como parar a divulgação do prazo antigo. Chamamos \textbf{LLMOps} ao trabalho de testar, lançar, observar e recuperar sistemas que integram modelos de linguagem. Inclui o prompt, a aplicação, os dados e as pessoas com autoridade para intervir.

Começa por guardar as versões que permitem reconstruir uma resposta. Um prompt num ficheiro versionado ajuda. Neste incidente, precisaremos também de saber qual era a política em vigor e qual foi usada pela aplicação.

\section{Ciclo de vida de um prompt em produção}

Um fluxo saudável (simplificado) para prompts críticos:

\begin{enumerate}[leftmargin=*]
\item \textbf{Design inicial}
Clarificas:
\begin{itemize}[leftmargin=*]
\item objetivo de negócio,
\item tipo de tarefa (classificação, extração, geração, etc.),
\item constraints (segurança, estilo, custos).
\end{itemize}
\item \textbf{Prototipagem}
Experimentas no chat:
\begin{itemize}[leftmargin=*]
\item crias alguns exemplos,
\item testas rapidamente variações,
\item recolhes casos fáceis e difíceis.
\end{itemize}
\item \textbf{Criação de golden set}
A partir dos exemplos, montas um conjunto de teste como no capítulo~\ref{ch:avaliacao}.
\item \textbf{Testes offline}
Corres:
\begin{itemize}[leftmargin=*]
\item prompt v1, v2, v3...
\item com um ou mais modelos,
\item medes métricas definidas (accuracy, taxas de erro, etc.).
\end{itemize}
\item \textbf{Revisão e aprovação}
Alguém além do autor do prompt revê:
\begin{itemize}[leftmargin=*]
\item comportamento,
\item riscos,
\item impacto esperado.
\end{itemize}
\item \textbf{Deploy controlado}
Lançamento em:
\begin{itemize}[leftmargin=*]
\item ambiente de teste / staging,
\item canary release (percentagem pequena de tráfego),
\item depois rollout total.
\end{itemize}
\item \textbf{Monitorização contínua}
Recolhes:
\begin{itemize}[leftmargin=*]
\item pedidos e respostas estritamente necessários, com campos sensíveis excluídos ou protegidos,
\item métricas de negócio,
\item feedback dos utilizadores.
\end{itemize}
\item \textbf{Iteração com disciplina}
Quando queres melhorar:
\begin{itemize}[leftmargin=*]
\item voltas ao golden set,
\item testas novas versões,
\item fazes rollout gradual,
\item manténs histórico.
\end{itemize}
\end{enumerate}

\section{Gestão de configuração e versionamento de prompts}

Prompts importantes não devem estar:

\begin{itemize}[leftmargin=*]
\item só numa UI de fornecedor,
\item apenas na cabeça de uma pessoa.
\end{itemize}

Boas práticas mínimas:

\begin{itemize}[leftmargin=*]
\item Guardar prompts em ficheiros de texto (por exemplo, \texttt{.txt}, \texttt{.md}, \texttt{.yaml}) no mesmo repositório que o código.
\item Dar nomes explícitos, por exemplo:
\begin{itemize}[leftmargin=*]
\item \texttt{router\_faq\_pt\_v3.txt},
\item \texttt{classificador\_pedidos\_suporte\_v5.yaml}.
\end{itemize}
\item Usar controlo de versão (Git):
\begin{itemize}[leftmargin=*]
\item ver diffs entre versões (o que mudou),
\item associar commits a tickets/tasks,
\item ter histórico de evolução.
\end{itemize}
\item Incluir metadados onde fizer sentido:
\begin{itemize}[leftmargin=*]
\item data,
\item autor,
\item modelo-alvo (por exemplo, \texttt{modelo\_e\_versao\_verificados}),
\item tarefas que o prompt destina-se a resolver,
\item links para golden sets ou testes.
\end{itemize}
\end{itemize}

Exemplo simplificado de ficheiro \texttt{.yaml}:

\begin{lstlisting}
id: classificador_tickets_suporte
version: 5
language: pt-PT
model_target: modelo_e_versao_verificados
owner: equipa_suporte
description: >
Classificação de tickets de suporte em categorias
operacionais para routing interno.

prompt_system: |
És um classificador de pedidos de suporte...
prompt_user_template: |
Classifica o texto abaixo...

tests_golden_set: data/golden_tickets_v2.json
\end{lstlisting}

\begin{checklistbox}{Antes de mexer num prompt de produção}
Garante que:

\begin{itemize}[leftmargin=*]
\item o prompt está em repositório versionado;
\item existe um golden set associado (nem que pequeno);
\item tens um ambiente para testar a nova versão antes de produção;
\item sabes como reverter rapidamente se algo correr mal.
\end{itemize}
\end{checklistbox}

\section{Registar o suficiente para reconstruir a resposta}
Para diagnosticar os dez dias, começamos pelo identificador do pedido, a versão do serviço e a da política usada. Juntamos a decisão de validação e o identificador da ferramenta, quando existe. Com esses campos, podemos seguir o percurso sem copiar toda a conversa para os logs.

A latência e os erros da API mostram como correu a execução. Para conhecer a qualidade, revê uma amostra com critérios e fontes. Uma queda nas escaladas pode significar que o serviço resolveu mais casos. Pode também significar que deixou de encaminhar casos difíceis. Confirma o significado com os resultados da revisão.

Define acesso e retenção para os registos. Truncar ou mascarar um texto pode reduzir o detalhe exposto. A avaliação de se os dados permitem identificar uma pessoa continua a depender do conteúdo e da possibilidade de os associar.

\section{Guardrails, políticas e fallbacks}

Mesmo com prompts excelentes, precisas de \emph{guardrails} fora do modelo:

\begin{itemize}[leftmargin=*]
\item filtros de conteúdo (antes e/ou depois do LLM),
\item whitelists/blacklists de operações e ferramentas,
\item mecanismos de fallback (para humanos, para respostas neutras).
\end{itemize}

Exemplos:

\begin{itemize}[leftmargin=*]
\item \textbf{Filtro de input}:
\begin{itemize}[leftmargin=*]
\item bloquear pedidos claramente fora de escopo (ex: tentar usar o chatbot de suporte para crime),
\item sanitizar partes do texto (ex: remover credenciais).
\end{itemize}
\item \textbf{Filtro de output}:
\begin{itemize}[leftmargin=*]
\item verificar se há padrões proibidos (ex: dados bancários, insultos),
\item moderar a resposta (ou substituí-la por algo neutro) se necessário.
\end{itemize}
\item \textbf{Fallbacks}:
\begin{itemize}[leftmargin=*]
\item se o modelo falha repetidamente num pedido, passar o caso a humano,
\item se há timeout, enviar mensagem clara ao utilizador.
\end{itemize}
\end{itemize}

\begin{checklistbox}{Guardrails básicos para um sistema com LLM}
Antes de colocar algo público:

\begin{itemize}[leftmargin=*]
\item tens filtros de input para bloquear conteúdo ilegal/abusivo?
\item tens filtros de output para evitar que o sistema devolva algo claramente inaceitável?
\item tens mecanismos de fallback para humanos e mensagens de erro decentes?
\item sabes quem é responsável por rever casos problemáticos?
\end{itemize}
\end{checklistbox}

\section{Ambientes, rollouts e gestão de risco}

Se aplicas boas práticas de DevOps, muito disto soa familiar:

\begin{itemize}[leftmargin=*]
\item \textbf{Ambientes separados}:
\begin{itemize}[leftmargin=*]
\item desenvolvimento,
\item staging (pré-produção),
\item produção.
\end{itemize}
\item \textbf{Canary releases}:
\begin{itemize}[leftmargin=*]
\item lançar nova versão de prompt/modelo para pequena percentagem de utilizadores,
\item observar métricas,
\item ampliar gradualmente se estiver estável,
\item reverter se algo correr mal.
\end{itemize}
\item \textbf{Feature flags}:
\begin{itemize}[leftmargin=*]
\item ativar/desativar funcionalidades baseadas em LLM sem redeploy completo,
\item experimentar novos fluxos com segmentos específicos de utilizadores.
\end{itemize}
\end{itemize}

\begin{keyideabox}{Trata “trocar de modelo” como um change de infra, não como detalhe}
Mudar de \emph{modelo} (por exemplo, de uma versão para outra) pode ter impacto tão grande quanto mudar de base de dados.
Planeia, testa, monitora e regista.
\end{keyideabox}

\section{Quem pode travar uma alteração}
Para cada fluxo, regista quem aprova uma versão e quem intervém num incidente. Na Lavandaria, a equipa técnica pode suspender o envio de respostas automáticas. A área que decide a política confirma o prazo em vigor. A equipa de apoio recebe os casos encaminhados e pode corrigir uma comunicação anterior.

Um lançamento gradual precisa de critérios de avanço e de paragem. Se o canário recuperar uma versão antiga do documento, suspende o fluxo afetado e confirma que a contenção funciona. Antes de reabrir, pede à mesma sonda que verifique a fonte em vigor. Um timeout numa ação externa exige confirmar o estado da operação antes de a repetir.

\section{O custo por tarefa concluída}
Calcula o custo de chamadas, ferramentas, infraestrutura e revisão humana por tarefa concluída dentro do critério. Um modelo barato que obriga a três tentativas e duas correções pode ser a opção mais cara. Compara candidatos nos mesmos casos e inclui as falhas. Na latência, acompanha a mediana, o p95 e a taxa de timeouts. O p95 ajuda a observar a cauda dos tempos mesmo quando a maioria das respostas é rápida.

Uma cache de respostas ou documentos deve incluir o âmbito de autorização, versão da fonte e prazo de validade. Se um cliente deixa de ter acesso a um documento, uma resposta antiga em cache não pode contornar essa revogação. Nunca uses uma cache partilhada como memória informal entre clientes.

Pseudonimizar reduz associação direta, mas os dados continuam pessoais quando podem ser reatribuídos. Regista por defeito identificadores técnicos e métricas; conserva texto bruto apenas quando necessário, com acesso, fundamento e retenção definidos. Para reproduzir incidentes, guarda uma amostra controlada, não um arquivo ilimitado de conversas~\citep{rgpd2016}.

\section{O incidente dos cinco dias}
\textbf{Ensaio executado.} A cache guarda a política fictícia \texttt{v1}, de dez dias úteis, com a chave \texttt{devolucoes}. Entra em vigor a \texttt{v2}, de cinco dias. Como a chave continua igual, a cache devolve a entrada antiga. A proposta tem todos os campos e o prazo é um inteiro positivo.

O \emph{canário} é um pedido de teste que compara documento, versão, SHA-256 e prazo com a política vigente, obtida fora da cache. Com a \texttt{v2} ativa, a sonda encontra a \texttt{v1} e o simulador encaminha o pedido sem devolver o prazo antigo.

\begin{table}[htbp]
\small
\begin{tabularx}{\textwidth}{@{}>{\raggedright\arraybackslash}p{.19\textwidth}>{\raggedright\arraybackslash}X>{\raggedright\arraybackslash}X>{\raggedright\arraybackslash}X@{}}
\toprule
Etapa & Versão vigente & Proposta & Decisão \\
\midrule
Base & v1, 10 dias úteis & v1, 10 dias úteis & Validada \\
Falha injetada & v2, 5 dias úteis & v1, 10 dias úteis & Encaminhada, sem prazo \\
Recuperação & v2, 5 dias úteis & v2, 5 dias úteis & Validada \\
\bottomrule
\end{tabularx}
\caption{Três etapas do ensaio sintético de cache. A proposta antiga está visível apenas no registo da sonda.}
\label{tab:incidente-cache}
\end{table}

A recuperação usa a chave composta pelo documento e pela versão. A entrada \texttt{v1} fica na cache, mas a consulta da \texttt{v2} recupera os cinco dias. Os hashes e os estados estão no registo do ensaio. Num serviço, reverter o código e mudar a política vigente são decisões separadas.

\begin{promptlabbox}{Reproduzir, conter e recuperar um incidente de cache}
\textbf{Objetivo e contexto:} reproduzir os três estados da tabela~\ref{tab:incidente-cache}. Requisitos: Python 3.10+ e biblioteca padrão. Os dados são sintéticos e a execução tem zero chamadas a modelos.

\textbf{Procedimento:} na raiz do projeto, executa o comando abaixo.
\par\noindent\texttt{python3 -m exemplos.laboratorio.operacao}\par
Compara a saída com \texttt{experiencias/operacao/resultado.json} e segue no código as chaves da cache.

\textbf{Resultado esperado e validação:} a proposta antiga passa a verificação estrutural e é encaminhada sem prazo. Depois da correção, o teste confirma versão \texttt{v2}, hash correspondente e cinco dias úteis. Experimenta propor dez dias com os identificadores da versão certa. A verificação do prazo deve travar também essa proposta.

\textbf{Falhas possíveis:} validar só o JSON, ler a versão de referência a partir da mesma cache ou restaurar o prazo antigo quando se reverte uma mudança de software.
\end{promptlabbox}
